\documentclass[12pt,a4paper]{exam}
\usepackage[utf8]{inputenc}
\usepackage{amsmath,amssymb,amsfonts}
\usepackage{geometry}
\usepackage{enumitem}
\usepackage{graphicx}
\usepackage{hyperref}
\usepackage{xcolor}
\geometry{a4paper, top=2cm, bottom=2cm, left=2.5cm, right=2.5cm}
\hypersetup{colorlinks=true, linkcolor=blue, urlcolor=blue}
\renewcommand{\thequestion}{\textbf{\arabic{question}}}
\renewcommand{\questionlabel}{\thequestion.}
\pointname{ marks}
\pointsdroppedatright
\bracketedpoints
\setlength{\parindent}{0pt}
\setlength{\emergencystretch}{3em}
\allowdisplaybreaks
\newsavebox{\schmathbox}
\newcommand{\fitmath}[1]{%
  \sbox{\schmathbox}{$\displaystyle #1$}%
  \ifdim\wd\schmathbox>\linewidth
    \resizebox{\linewidth}{!}{\usebox{\schmathbox}}%
  \else
    \usebox{\schmathbox}%
  \fi}

\begin{document}

\thispagestyle{empty}
\begin{center}
  \textbf{\LARGE Diag} \\[0.3cm]
  \textbf{Time Allowed: 3 Hours} \hfill \textbf{Maximum Marks: 80} \\[0.3cm]
  \rule{\textwidth}{0.5pt}
\end{center}

\vspace{0.3cm}

\noindent\textbf{General Instructions:}
\begin{enumerate}
  \item {\normalfont\normalcolor This question paper contains 40 questions. All questions are compulsory.}
  \item {\normalfont\normalcolor Questions are grouped into sections A through E.}
  \item {\normalfont\normalcolor Use of calculators is NOT permitted.}
\end{enumerate}
\bigskip
\noindent\textbf{\large Section A: Multiple Choice \& Assertion-Reason (1 mark each)}

\begin{questions}
\question[1] The HCF of the smallest prime number and the smallest composite number is:
\begin{enumerate}[label=(\Alph*)]
  \item (A) 2
  \item (B) 1
  \item (C) 4
  \item (D) 8
\end{enumerate}
\vspace{0.15cm}

\question[1] If two positive integers $a$ and $b$ are written as $a = x^3y^2$ and $b = xy^3$, where $x, y$ are prime numbers, then the $\text{HCF}(a, b)$ is:
\begin{enumerate}[label=(\Alph*)]
  \item (A) xy
  \item (B) xy\textasciicircum{}3
  \item (C) xy\textasciicircum{}2
  \item (D) x\textasciicircum{}2y\textasciicircum{}2
\end{enumerate}
\vspace{0.15cm}

\question[1] The decimal expansion of the rational number $\frac{13}{125}$ will terminate after how many decimal places?
\begin{enumerate}[label=(\Alph*)]
  \item (A) 2
  \item (B) 3
  \item (C) 4
  \item (D) 1
\end{enumerate}
\vspace{0.15cm}

\question[1] If the $\text{LCM}(26, 169) = 338$, then $\text{HCF}(26, 169)$ is:
\begin{enumerate}[label=(\Alph*)]
  \item (A) 26
  \item (B) 13
  \item (C) 39
  \item (D) 1
\end{enumerate}
\vspace{0.15cm}

\question[1] The product of a non-zero rational and an irrational number is always:
\begin{enumerate}[label=(\Alph*)]
  \item (A) rational
  \item (B) integer
  \item (C) irrational
  \item (D) one
\end{enumerate}
\vspace{0.15cm}

\question[1] Assertion (A): If $p(x) = x^3 - 3x^2 + 5x - 3$ is divided by $x-1$, the remainder is 0. Reason (R): If $p(a) = 0$, then $(x-a)$ is a factor of $p(x)$.
\begin{enumerate}[label=(\Alph*)]
  \item (A) Both Assertion (A) and Reason (R) are true and Reason (R) is the correct explanation of Assertion (A)
  \item (B) Both Assertion (A) and Reason (R) are true but Reason (R) is NOT the correct explanation of Assertion (A)
  \item (C) Assertion (A) is true but Reason (R) is false
  \item (D) Assertion (A) is false but Reason (R) is true
\end{enumerate}
\vspace{0.15cm}

\question[1] Assertion (A): If a polynomial $p(x)$ is divided by a linear polynomial $g(x)$, the remainder is a constant. Reason (R): The degree of a linear polynomial is 1.
\begin{enumerate}[label=(\Alph*)]
  \item (A) Both Assertion (A) and Reason (R) are true and Reason (R) is the correct explanation of Assertion (A)
  \item (B) Both Assertion (A) and Reason (R) are true but Reason (R) is NOT the correct explanation of Assertion (A)
  \item (C) Assertion (A) is true but Reason (R) is false
  \item (D) Assertion (A) is false but Reason (R) is true
\end{enumerate}
\vspace{0.15cm}

\question[1] Assertion (A): The HCF of $12$ and $18$ is $6$. Reason (R): For any two positive integers $a$ and $b$, $HCF(a, b) \times LCM(a, b) = a \times b$.
\begin{enumerate}[label=(\Alph*)]
  \item (A) Both Assertion (A) and Reason (R) are true and Reason (R) is the correct explanation of Assertion (A)
  \item (B) Both Assertion (A) and Reason (R) are true but Reason (R) is NOT the correct explanation of Assertion (A)
  \item (C) Assertion (A) is true but Reason (R) is false
  \item (D) Assertion (A) is false but Reason (R) is true
\end{enumerate}
\vspace{0.15cm}

\question[1] Assertion (A): $\sqrt{3}$ is an irrational number. Reason (R): The square root of any positive integer is always an irrational number.
\begin{enumerate}[label=(\Alph*)]
  \item (A) Both Assertion (A) and Reason (R) are true and Reason (R) is the correct explanation of Assertion (A)
  \item (B) Both Assertion (A) and Reason (R) are true but Reason (R) is NOT the correct explanation of Assertion (A)
  \item (C) Assertion (A) is true but Reason (R) is false
  \item (D) Assertion (A) is false but Reason (R) is true
\end{enumerate}
\vspace{0.15cm}

\question[1] Assertion (A): The number $0.375$ is a terminating decimal. Reason (R): A rational number $p/q$ terminates if the prime factorization of $q$ is of the form $2^n 5^m$, where $n, m$ are non-negative integers.
\begin{enumerate}[label=(\Alph*)]
  \item (A) Both Assertion (A) and Reason (R) are true and Reason (R) is the correct explanation of Assertion (A)
  \item (B) Both Assertion (A) and Reason (R) are true but Reason (R) is NOT the correct explanation of Assertion (A)
  \item (C) Assertion (A) is true but Reason (R) is false
  \item (D) Assertion (A) is false but Reason (R) is true
\end{enumerate}
\vspace{0.15cm}

\question[1] Assertion (A): The HCF of two consecutive natural numbers is 1. Reason (R): Any two consecutive natural numbers are co-prime.
\begin{enumerate}[label=(\Alph*)]
  \item (A) Both Assertion (A) and Reason (R) are true and Reason (R) is the correct explanation of Assertion (A)
  \item (B) Both Assertion (A) and Reason (R) are true but Reason (R) is NOT the correct explanation of Assertion (A)
  \item (C) Assertion (A) is true but Reason (R) is false
  \item (D) Assertion (A) is false but Reason (R) is true
\end{enumerate}
\vspace{0.15cm}

\question[1] Assertion (A): The product of a non-zero rational and an irrational number is always irrational. Reason (R): $\sqrt{2}$ is an irrational number.
\begin{enumerate}[label=(\Alph*)]
  \item (A) Both Assertion (A) and Reason (R) are true and Reason (R) is the correct explanation of Assertion (A)
  \item (B) Both Assertion (A) and Reason (R) are true but Reason (R) is NOT the correct explanation of Assertion (A)
  \item (C) Assertion (A) is true but Reason (R) is false
  \item (D) Assertion (A) is false but Reason (R) is true
\end{enumerate}
\vspace{0.15cm}

\end{questions}
\bigskip
\noindent\textbf{\large Section B: Very Short Answer (2 marks each)}

\begin{questions}
\question[2] On dividing the polynomial $p(x) = x^3 - 3x^2 + x + 2$ by a polynomial $g(x)$, the quotient and remainder were $(x - 2)$ and $(-2x + 4)$ respectively. Find $g(x)$.
\vspace{0.15cm}

\question[2] Check whether $g(x) = x^2 + 3x + 1$ is a factor of the polynomial $p(x) = 3x^4 + 5x^3 - 7x^2 + 2x + 2$ by applying the division algorithm.
\vspace{0.15cm}

\question[2] If the polynomial $f(x) = x^4 - 6x^3 + 16x^2 - 25x + 10$ is divided by another polynomial $x^2 - 2x + k$, the remainder comes out to be $x + a$. Find the values of $k$ and $a$.
\vspace{0.15cm}

\question[2] Find the quotient and remainder when $p(x) = x^3 - x^2 + x - 1$ is divided by $g(x) = x - 1$.
\vspace{0.15cm}

\question[2] What must be subtracted from $p(x) = 8x^4 + 14x^3 - 2x^2 + 7x - 8$ so that the resulting polynomial is exactly divisible by $g(x) = 4x^2 + x - 2$?
\vspace{0.15cm}

\question[2] Find the HCF of $1260$ and $7344$ using the Euclid's division algorithm.
\vspace{0.15cm}

\question[2] Show that any positive odd integer is of the form $4q + 1$ or $4q + 3$, where $q$ is some integer.
\vspace{0.15cm}

\question[2] Check whether $12^n$ can end with the digit $0$ for any natural number $n$.
\vspace{0.15cm}

\question[2] Given that $\text{HCF}(306, 657) = 9$, find $\text{LCM}(306, 657)$.
\vspace{0.15cm}

\question[2] Without performing long division, state whether $\frac{17}{8}$ has a terminating or non-terminating repeating decimal expansion.
\vspace{0.15cm}

\end{questions}
\bigskip
\noindent\textbf{\large Section C: Short Answer (3 marks each)}

\begin{questions}
\question[3] On dividing $x^3 - 3x^2 + x + 2$ by a polynomial $g(x)$, the quotient and remainder were $(x - 2)$ and $(-2x + 4)$ respectively. Find $g(x)$.
\vspace{0.15cm}

\question[3] Find all the zeroes of the polynomial $p(x) = 2x^4 - 3x^3 - 3x^2 + 6x - 2$, if you know that two of its zeroes are $\sqrt{2}$ and $-\sqrt{2}$.
\vspace{0.15cm}

\question[3] What must be subtracted from $f(x) = x^4 + 2x^3 - 13x^2 - 12x + 21$ so that the resulting polynomial is exactly divisible by $g(x) = x^2 - 4x + 3$?
\vspace{0.15cm}

\question[3] If the polynomial $x^4 + x^3 + 8x^2 + ax + b$ is divisible by $x^2 + 1$, find the values of $a$ and $b$.
\vspace{0.15cm}

\question[3] Check whether $g(t) = t^2 - 3$ is a factor of $p(t) = 2t^4 + 3t^3 - 2t^2 - 9t - 12$ by applying the division algorithm.
\vspace{0.15cm}

\question[3] Prove that $3 + 2\sqrt{5}$ is an irrational number, given that $\sqrt{5}$ is an irrational number.
\vspace{0.15cm}

\question[3] Find the HCF and LCM of 96 and 404 by the prime factorization method and verify that HCF $\times$ LCM = product of the two numbers.
\vspace{0.15cm}

\question[3] Without performing long division, state whether $\frac{129}{2^2 \times 5^7 \times 7^5}$ has a terminating or non-terminating repeating decimal expansion. Justify your answer.
\vspace{0.15cm}

\question[3] Use Euclid's division lemma to show that the square of any positive integer is either of the form $3m$ or $3m + 1$ for some integer $m$.
\vspace{0.15cm}

\question[3] There is a circular path around a sports field. Sonia takes 18 minutes to drive one round of the field, while Ravi takes 12 minutes for the same. If they both start at the same point and at the same time, and go in the same direction, after how many minutes will they meet again at the starting point?
\vspace{0.15cm}

\end{questions}
\bigskip
\noindent\textbf{\large Section D: Long Answer (4-5 marks each)}

\begin{questions}
\question[5] Obtain all the zeroes of the polynomial $p(x) = 2x^4 - 3x^3 - 3x^2 + 6x - 2$, if you know that two of its zeroes are $\sqrt{2}$ and $-\sqrt{2}$.
\vspace{0.15cm}

\question[5] On dividing $x^3 - 3x^2 + x + 2$ by a polynomial $g(x)$, the quotient and remainder were $(x - 2)$ and $(-2x + 4)$ respectively. Find $g(x)$.
\vspace{0.15cm}

\question[5] Find the values of $a$ and $b$ so that the polynomial $x^4 + x^3 + 8x^2 + ax + b$ is exactly divisible by $x^2 + 1$.
\vspace{0.15cm}

\question[5] Given that $x^2 + 1$ is a factor of $x^4 - 2x^3 + x^2 - 2x + 2$, use the division algorithm to find what must be subtracted from the polynomial so that the result is exactly divisible by $x^2 + 1$.
\vspace{0.15cm}

\question[5] Prove that $\sqrt{5}$ is an irrational number. Hence, show that $3 + 2\sqrt{5}$ is also irrational.
\vspace{0.15cm}

\question[5] Two tankers contain $850$ litres and $680$ litres of petrol respectively. Find the maximum capacity of a container which can measure the petrol of either tanker in exact number of times.
\vspace{0.15cm}

\question[5] Show that any positive odd integer is of the form $4q + 1$ or $4q + 3$, where $q$ is some integer.
\vspace{0.15cm}

\question[5] Find the smallest number which when divided by 28 and 32 leaves remainders 8 and 12 respectively.
\vspace{0.15cm}

\end{questions}
\bigskip

\end{document}
